% update_conclusion.tex — Выводы по секции 3.1 (INTENSE)

\paragraph{Анализ индивидуальной нейронной селективности (секция~\ref{sec:intense}).}

Разработан информационно-теоретический фреймворк INTENSE для выявления и количественной оценки нейронной селективности в данных кальциевого имиджинга при свободном поведении. Основные результаты:

\begin{enumerate}
    \item Предложен метод оценки нейронной селективности на основе взаимной информации с использованием копула-энтропии, работающий непосредственно с сырым кальциевым сигналом без необходимости деконволюции спайков и позволяющий обрабатывать переменные различных типов (непрерывные, дискретные, многомерные) в рамках единой аналитической схемы.

    \item Разработана двухэтапная процедура статистического тестирования с пермутациями на основе циклического сдвига, сохраняющими автокорреляцию нейронных и поведенческих сигналов, ускоренная методом быстрого преобразования Фурье в $100$--$2500$ раз.

    \item Предложена процедура разделения истинной смешанной селективности и артефактов, обусловленных корреляциями поведенческих переменных, на основе условной взаимной информации и информации взаимодействия.

    \item При применении к гиппокампальным записям 16 мышей при свободном исследовании среды были выявлены разнообразные типы селективных нейронов (клетки места, клетки направления головы, нейроны, модулируемые скоростью, объектные клетки). Показано, что примерно треть наблюдаемой мультиселективности объясняется поведенческой избыточностью, а замена взаимной информации на линейную корреляцию приводит к потере более двух третей значимых ассоциаций.

    \item Синтетическая валидация в широком диапазоне отношений сигнал/шум и надёжности ответа подтвердила робастность метода и критическую важность контроля временной автокорреляции при тестировании значимости.
\end{enumerate}
